\section{The unit-zeta cubic locus consists of boundary points}
\label{cb19-section}

This section applies the complete fixed-curve classification of
Section~\ref{sec-fixed-rational-locus} to the original, common-source
biquadratic curve. The definitions and projective inverses are those
of Sections~\ref{rl-section}, \ref{hg-section}, \ref{eqg-section}
and~\ref{ssg-section}. Put \(\zeta^2-\zeta+1=0\), with the same
Eisenstein coordinate convention as those sections.

\begin{proposition}[The six points in the original first quotient]
\label{cb-first-quotient}
For the smooth projective curve
\[
 H:y^2=P(s):=s^6+3s^5-3s^4-11s^3-3s^2+3s+1,
\]
the complete rational locus consists of its two points over each of
\(s=0,-1,\infty\).
\end{proposition}
\begin{proof}
Writing \(S_C(z)=z^6-27z^4+99z^2-9\), the exact dense-chart maps are
\begin{equation}
 z=\frac{s-1}{s+1},\quad W=\frac{8y}{(s+1)^3},\qquad
 s=\frac{1+z}{1-z},\quad y=\frac W{(1-z)^3}.
 \label{cb-chart}
\end{equation}
They are the actual transformation of Theorem~\ref{eqg-quotients}:
its intermediate coordinates are
\(t=(s+1)/(s-1)\), \(v=y(t-1)^3\), with
\(v^2=-9t^6+99t^4-27t^2+1\). Substitution \(z=1/t\),
\(W=v/t^3\) gives \eqref{cb-chart}. Equivalently,
\[
 (s+1)^6S_C\left(\frac{s-1}{s+1}\right)=64P(s).
\]
The inverse dense maps extend to an isomorphism of smooth projective
curves, by \cite[Lemmas~53.2.2 and 53.2.6]{StacksCurves2026}.

The complete list
\(C(\mathbb Q)=\{(\pm1,\pm8),\infty_+,\infty_-\}\)
therefore transfers as follows:
\[
\begin{array}{c|c}
 \text{point of }C&\text{point of }H\\\hline
 z=-1,\ W=\pm8&s=0,\ y=\pm1\\
 z=1,\ W=\pm8&s=\infty,\ y/s^3=\pm1\\
 z=\infty,\ W/z^3=\epsilon&s=-1,\ y=-\epsilon
\end{array}
\]
The exceptional rows follow from the actual chart limits
\(y/s^3=W/(1+z)^3\) and
\(y=(W/z^3)/(1/z-1)^3\). Thus every apparent zero denominator is
accounted for, and the six listed points exhaust \(H(\mathbb Q)\).
\end{proof}

\begin{theorem}[The complete common-source cubic locus]
\label{cb-twelve-points}
Let
\begin{align}
 A_0&=S^3-3ST^2-T^3,&B_0&=3ST(S+T),\\
 A&=-B_0,&B&=A_0+B_0,&r&=A/B.
 \label{cb-forms}
\end{align}
The smooth projective normalization \(D=D_{3,\zeta}\) of
\begin{equation}
 v^2=1+r,\qquad w^2=1-3r
 \label{cb-biquadratic}
\end{equation}
has exactly twelve rational points. Their source is
\([S:T]=[0:1],[-1:1],[1:0]\), and at each source
\((v,w)\) takes all four values in \(\{1,-1\}^2\).
Their original projective parameter \(t\) is given by
\begin{equation}
\begin{array}{cc|c}
 v&w&t\\\hline
 1&1&\infty\\1&-1&1\\-1&1&-1\\-1&-1&-1/2
\end{array}
 \label{cb-boundary-table}
\end{equation}
\end{theorem}
\begin{proof}
The first quotient of \eqref{cb-biquadratic} is exactly \(H\),
with \(y_1=Bv\) and \(y_1^2=A_0(A_0+B_0)\).
At the same source, the second quotient has \(y_2=Bw\) and
\begin{equation}
 y_2^2=(A_0+B_0)(A_0+4B_0).
 \label{cb-second-square}
\end{equation}
Every rational point of \(D\) maps to a rational point of \(H\)
under its projective quotient morphism. Proposition~\ref{cb-first-quotient}
therefore restricts its source to the following three values:
\[
\begin{array}{c|rrrr}
 s&A_0&B_0&A&B\\\hline
 0&-1&0&0&-1\\-1&1&0&0&1\\\infty&1&0&0&1
\end{array}
\]
In particular \(B\ne0\) and \(r=0\) at each source. The second
square \eqref{cb-second-square} is then \(B^2\), giving both
rational choices \(w=\pm1\); the two first-quotient points give
\(v=\pm1\). Points from different sources have not been paired.
These are four distinct rational points at each source. Both square
roots are nonzero, and \(r=0\) avoids the branch values
\(-1,1/3,\infty\). The raw equations are smooth there, so normalization
neither merges these points nor introduces another local branch.
This proves the exact total of twelve.

For the two middle rows of \eqref{cb-boundary-table}, the original
inverse \eqref{hg-inverse} is
\(t=(v+w+2)/(v-w)\). The first and last rows require its genuine
projective continuation; the last gives an apparent \(0/0\).
Use the original formulas instead:
\[
 v=\frac{t^2+2t}{t^2+t+1},\qquad
 w=\frac{t^2-2t-2}{t^2+t+1}.
\]
At infinity both ratios are one, and at \(t=-1/2\) both are minus
one. Also
\[
 r(t)=\frac{(t^2-1)(2t+1)}{(t^2+t+1)^2}
\]
has precisely the four simple projective zeros
\(\infty,1,-1,-1/2\). These give all projective inverse values,
including the removable \(0/0\) case.
\end{proof}

\begin{corollary}[Exclusion of the actual positive unit-zeta branch]
\label{cb-positive-empty}
The curve \(D_{3,\zeta}\), and its Gaussian cover over the same
positive domain, have no rational point with finite original
parameter \(t>1\). In particular there are no positive coprime
integers \(a,b\) such that
\[
 M=a^2+ab+b^2=U^2,\qquad F(a,b)=Q^3,\qquad Q>1,
\]
whose actual oriented second Eisenstein decomposition has unit class
\(\zeta\) modulo cubes of units. More generally the same positive
locus of \(D_{g,\zeta}\) is empty for every positive odd integer
\(g\) divisible by three.
\end{corollary}
\begin{proof}
The exact positive inverse of Theorem~\ref{hg-positive} requires
finite \(t>1\), equivalently \(0<r<1/3\) with \(v>0\).
None of the points in \eqref{cb-boundary-table} has this property.
For clarity, the homogeneous first-root output is
\[
 a_0=T_0^2-W_0^2,\qquad b_0=W_0(2T_0+W_0).
\]
At the four values of \(t=T_0/W_0\) in the table, its primitive
signed pairs are respectively
\((1,0),(0,1),(0,-1),(-1,0)\). The finite representatives at
\(t=1,-1/2\) have content three, which must be divided out.
Every primitive boundary pair has \(ab=0\) and \(M=F=1\);
none is a positive seed or has \(Q>1\).

An actual seed in the asserted unit class has, after absorbing a
minus sign into the root,
\[
 ab+M\zeta=\zeta w_1^3.
\]
The exact projective inverse of Section~\ref{rl-section} then
supplies the common source and \(t=(a+U)/b>1\), contradicting
Theorem~\ref{cb-twelve-points}. That inverse included primitive,
ramified, unit and infinite root inputs before imposing the positive
domain. Thus none of these input exceptions restores a positive
point. The previously proved converse on that domain makes this an
exclusion of the actual branch, not just a necessary auxiliary square.

For \(g=3k\) odd, the projective power maps satisfy the exact
identity
\[
 L_{\zeta,g}=L_{\zeta,3}\circ L_{1,k}.
\]
It gives a nonconstant rational curve map
\((t,s)\mapsto(t,L_{1,k}(s))\) from \(D_{g,\zeta}\) to
\(D_{3,\zeta}\), extending to the smooth projective models as in
\cite[Lemmas~53.2.2 and 53.2.6]{StacksCurves2026}.
It retains \(t\); a rational point with finite \(t>1\) would map
to the excluded positive locus. No enumeration of the boundary
preimages for larger \(g\) is needed or asserted.
\end{proof}

\paragraph{Scope.}
The identity cubic unit class already has its separate complete
three-adic obstruction. The distinct class \(\zeta^2\) is not
eliminated here: conjugation changes the target ratio \(r\) to
\(-1-r\), which does not preserve \((0,1/3)\). It cannot silently
be substituted for a positive-domain isomorphism. General second
residuals, the ramified first norm, other exponents and uniform
point-height bounds remain outside this conclusion. The theorem
uses the complete ordinary and exact finite fixed-curve proof,
together with the already proved projective inverses; it is not a
new Lean formalization of a curve, rational locus or ABC statement.
